% !TEX root = ../main.tex
% =====================================================================
% CAPITOLO 6 — ORION: IMPLEMENTAZIONE
% =====================================================================

\chapter{Orion: implementazione}
\label{chap:orion-implementazione}

Questo capitolo descrive come le scelte architetturali presentate nel Capitolo~\ref{chap:orion-architettura} 
sono state tradotte in codice.

% =====================================================================
\section{Organizzazione multi-repository}
\label{sec:orion-impl-multirepo}
% =====================================================================

Con l'aggiunta di Orion, viene completato il disegno obiettivo del progetto.
Ai repository descritti nel Capitolo~\ref{chap:agenti-implementazione} (Fig.~\ref{fig:multirepo_tree})
si aggiungono quattro nuovi servizi dedicati all'orchestratore:

\begin{figure}[h]
  \centering
  \begin{lstlisting}[language=bash]
vera-backend/              # Vera Backend (FastAPI + ADK + Gemini)
vera-frontend/             # Vera Frontend (React)
vera-admin-backend/        # Pannello amministrativo Backend
vera-admin-frontend/       # Pannello amministrativo Frontend
sysaid-backend/            # SysAid Backend unificato (Gemini + Sonnet)
sysaid-frontend/           # SysAid Frontend
coopolicy-backend/         # CooPolicy Backend (FastAPI + RAG)
coopolicy-frontend/        # CooPolicy Frontend
allycare-backend/          # AllyCare Backend (FastAPI + ADK + BigQuery)
allycare-frontend/         # AllyCare Frontend
portal-frontend/           # Portale utenti (MSAL)
admin-portal-frontend/     # Portale amministrativo
orion-backend/             # Orchestratore Orion Backend
orion-frontend/            # Orchestratore Orion Frontend
orion-admin-backend/       # Orion Admin Panel Backend
orion-admin-frontend/      # Orion Admin Panel Frontend
infrastructure/            # Terraform per infrastruttura
  \end{lstlisting}
  \caption[Elenco completo dei repository della suite]{Elenco completo dei repository della suite a seguito dell'aggiunta di Orion. 
  Ogni servizio è un progetto Git autonomo con proprio ciclo di CI/CD e branch per ambiente (\texttt{test}, \texttt{main}).}
  \label{fig:multirepo_tree_completo}
\end{figure}

% =====================================================================
\section{Backend Orion}
\label{sec:orion-impl-backend}
% =====================================================================

Il backend di Orion è stato sviluppato a partire dal backend di Vera;
rispetto a Vera, il backend aggiunge tre componenti nuovi: il \emph{router agent},
il \emph{proxy} verso i sub-agenti e l'endpoint dedicato alla valutazione.
Tutti gli endpoint utente sono esposti sotto il prefisso \texttt{/api/orion/};
gli endpoint amministrativi sono serviti da un servizio Cloud Run dedicato (\texttt{orion-admin-backend}),
descritto nella Sez.~\ref{sec:orion-impl-admin}.

\subsection{Router agent}
\label{subsec:orion-impl-router}

Il modulo \texttt{router\_agent.py} implementa la funzione asincrona \texttt{classify}, 
che costituisce il cuore del routing step. La funzione riceve tre parametri: il prompt dell'utente, la cronologia della conversazione (ultimi sei turni) 
e la lista degli agenti disponibili.

Il processo di classificazione si articola in quattro fasi:

\begin{enumerate}
  \item \textbf{Costruzione del prompt.} Il sistema compone il prompt completo 
  concatenando il \emph{system prompt} del router (che descrive le regole di classificazione, 
  la semantica della confidence e le istruzioni per il reasoning), 
  il blocco descrittivo degli agenti disponibili (nome, descrizione, esempi di richieste tipiche per ciascuno) 
  e il prompt dell'utente. La cronologia della conversazione è inclusa in forma sintetica 
  (ultimi sei turni, ciascuno troncato a 200~caratteri, configurabile) per consentire al modello di risolvere 
  riferimenti incrociati nei turni di \emph{follow-up}.

  \item \textbf{Chiamata a Vertex AI con output strutturato.} La funzione invoca il modello Gemini~2.5 Flash~\cite{team2024gemini} 
  attraverso il client \texttt{google.genai}, configurato con \texttt{temperature = 0.1} (più deterministico), \texttt{max\_output\_tokens = 1024} e, 
  soprattutto, \texttt{response\_mime\_type = "application/json"} con un \texttt{response\_schema} esplicito che vincola il modello a produrre un JSON 
  conforme allo schema descritto nella Sez.~\ref{subsec:orion-output-strutturato}. 
  Il campo \texttt{agent} è vincolato a un'\texttt{enum} contenente i soli agenti passati come parametro, eliminando la possibilità di target inventati.

  \item \textbf{Parsing e validazione.} La risposta JSON viene parsificata e 
  validata attraverso il modello Pydantic \texttt{RouterOutput}. 
  Se l'agente restituito non è nella lista degli agenti disponibili (possibile in caso di allucinazione residua nonostante lo schema), 
  il sistema effettua un \emph{fallback} su \texttt{"self"}. Se la confidence è inferiore alla soglia configurata e il modello non ha proposto una domanda di chiarimento, 
  il sistema ne genera una automaticamente.

  \item \textbf{Logging.} La decisione di routing viene loggata con livello \texttt{INFO} 
  se la confidence è $\geq 0{,}80$, \texttt{WARNING} altrimenti. Il log include agente selezionato, complessità, 
  confidence e flag di clarification, e alimenta la pagina \emph{Routing Logs} del pannello di amministrazione.
\end{enumerate}

In caso di eccezione durante la classificazione (errore di rete, timeout verso Vertex AI, errore di parsing), 
il modulo restituisce un \emph{fallback deterministico}: \texttt{agent = "self"}, \texttt{confidence = 1.0}, 
con un reasoning che documenta l'errore.

\begin{lstlisting}[language=Python,caption={Firma della funzione \texttt{classify} e struttura dell'output.},label={lst:classify-signature}]
class RouterOutput(BaseModel):
    agent: str        # "self" | "sysaid" | "allycare" | "coopolicy"
    complexity: str   # "low" | "high"
    confidence: float # 0.0 .. 1.0
    reasoning: str    # Chain-of-Thought esplicita
    clarification_needed: Optional[str] = None

async def classify(
    prompt: str,
    history: list,
    available_agents: list[str],
) -> RouterOutput:
\end{lstlisting}

\subsection{Agent proxy}
\label{subsec:orion-impl-proxy}

Il modulo \texttt{agent\_proxy.py} implementa la funzione \texttt{forward}, 
che inoltra la richiesta al sub-agente selezionato dal router e normalizza la risposta. 
La tabella degli endpoint è configurata tramite variabili d'ambiente, una per ciascun sub-agente:

\begin{table}[htbp]
  \centering
  \caption[Mapping degli endpoint dei sub-agenti]{Mapping tra agente target, variabile d'ambiente per l'URL del backend e path dell'endpoint.}
  \label{tab:agent-endpoints}
  \begin{tabular}{lll}
    \toprule
    \textbf{Agent key} & \textbf{Env var}             & \textbf{Path} \\
    \midrule
    \texttt{sysaid}    & \texttt{SYSAID\_BACKEND\_URL} & \texttt{/api/sysaid/prompt} \\
    \texttt{allycare}  & \texttt{NPS\_BACKEND\_URL}    & \texttt{/api/nps/prompt} \\
    \texttt{coopolicy} & \texttt{DOCS\_BACKEND\_URL}   & \texttt{/api/docs/prompt} \\
    \bottomrule
  \end{tabular}
\end{table}

Il payload inviato al sub-agente contiene, oltre al prompt e alla cronologia, 
i campi \texttt{complexity}, \texttt{user\_name}, \texttt{user\_role} e \texttt{user\_direction}, 
consentendo al sub-agente di personalizzare la risposta senza dover decodificare un token JWT.

La funzione di normalizzazione \texttt{\_normalize\_response} mappa le \emph{response shape} 
eterogenee dei sub-agenti verso un oggetto \texttt{OrionProxyResponse} con struttura fissa. 
I campi opzionali (\texttt{queries}, \texttt{charts}, \texttt{sources}, \texttt{data}) 
sono valorizzati solo se presenti nella risposta del sub-agente, 
consentendo al frontend di adattare il rendering in base all'agente che ha prodotto la risposta.

\subsection{Integrazione nel flusso di chat}
\label{subsec:orion-impl-chat}

L'endpoint SSE \texttt{/api/orion/chat/stream} è il punto di ingresso per le richieste dell'utente. 
Il flusso, rispetto a Vera, aggiunge un passo di routing dopo i controlli di quota e di sicurezza e prima della generazione:

\begin{enumerate}
  \item Verifica quota token dell'utente.
  \item Sanitizzazione del prompt (Model Armor).
  \item \textbf{[Nuovo]} Chiamata a \texttt{router\_agent.classify()}.
  \item \textbf{[Nuovo]} Logging della decisione in \texttt{orion\_routing\_logs}.
  \item Se \texttt{clarification\_needed}: ritorno della domanda all'utente, fine del turno.
  \item Se \texttt{agent != "self"}: chiamata a \texttt{agent\_proxy.forward()}, ritorno della risposta normalizzata.
  \item Se \texttt{agent == "self"} e la conversazione ha documenti caricati: pipeline RAG interna (Gemini Context Caching).
  \item Se \texttt{agent == "self"} senza documenti: pipeline generalista (ricerca web, conoscenza parametrica, streaming SSE).
\end{enumerate}

La Fig.~\ref{fig:security_chain_orion} visualizza il flusso completo come evoluzione della 
catena di sicurezza di Vera (Fig.~\ref{fig:security_chain}): il \emph{routing step} si innesta fra il 
controllo di sicurezza sul prompt e la generazione. 
Le tre diramazioni del router riconfluiscono prima della sanitizzazione dell'output, così che ogni risposta, 
inclusa quella prodotta da un sub-agente, attraversi i medesimi controlli post-generazione e lo stesso audit della pipeline originaria.

\begin{figure}[htbp]
    \centering
    \includegraphics[width=0.95\textwidth,height=0.9\textheight,keepaspectratio]{flusso-prompt-full-orion}
    \caption[Catena di sicurezza e routing di Orion]{Il flusso completo di una richiesta in Orion: verifica della quota, sanitizzazione del prompt tramite Model Armor 
    (RAI, prompt injection, DLP) con esenzioni per gruppo, \emph{routing step} che instrada verso la pipeline interna, un sub-agente specializzato oppure una richiesta di chiarimento, 
    e infine sanitizzazione dell'output e audit su \texttt{chat\_logs}. Evoluzione della Fig.~\ref{fig:security_chain} con l'aggiunta dello strato di orchestrazione.}
    \label{fig:security_chain_orion}
\end{figure}

Il routing step viene eseguito \emph{sempre}, anche in presenza di documenti caricati: 
ciò consente all'utente di porre, nella stessa conversazione, sia domande sui propri documenti sia richieste destinate a un sub-agente. 
La selezione tra la pipeline RAG interna e la pipeline generalista avviene \emph{dopo} il routing, sulla base del campo \texttt{agent} restituito dal router.

\subsection{Endpoint per la valutazione}
\label{subsec:orion-impl-classify}

L'endpoint \texttt{POST /api/orion/classify} espone il solo routing step, senza eseguire la risposta dell'agente selezionato. 
Riceve un prompt e una storia opzionale, e restituisce l'output strutturato del router. 
L'endpoint è stato progettato specificamente per il framework di valutazione descritto nel Capitolo~\ref{chap:orion-valutazione}: 
testare il router in isolamento consente di eseguire centinaia di test case con il costo di una sola chiamata Flash per test case, senza attivare i sub-agenti.

% =====================================================================
\section{Frontend Orion}
\label{sec:orion-impl-frontend}
% =====================================================================

Il frontend di Orion è un'applicazione React~19 con TypeScript, derivata dal frontend di Vera e 
adattata per la nuova architettura. Le differenze significative riguardano tre aspetti: 
il rendering del badge di routing, la gestione delle \emph{response shape} eterogenee e il supporto alle citazioni di CooPolicy.

\subsection{Badge di routing e trasparenza}
\label{subsec:orion-impl-fe-routing}

Ogni messaggio di risposta include un \emph{badge} visibile sopra il testo che 
indica quale agente ha servito la richiesta e con quale livello di complessità: 
ad esempio, ``SysAid $\cdot$ Alta'' o ``CooPolicy $\cdot$ Bassa''. 
Il badge è un componente React che riceve i campi \texttt{agent\_used} e \texttt{complexity} 
dall'oggetto \texttt{OrionProxyResponse}. La scelta di rendere visibile il routing è motivata da due ragioni: il valore dimostrativo nel contesto della tesi, 
e la fiducia dell'utente nel sistema (l'utente vede che la risposta proviene da un agente specializzato).

\begin{figure}[htbp]
  \centering
  \includegraphics[width=0.7\textwidth,height=0.9\textheight,keepaspectratio]{routing1}
  \caption[Orion: cambio agente nella stessa conversazione]{Routing trasparente con cambio di agente all'interno di un'unica conversazione: la prima richiesta (NPS) è instradata ad AllyCare (badge \emph{ORION $\cdot$ AllyCare $\cdot$ Bassa}), la successiva (ticket IT) a SysAid (badge \emph{ORION $\cdot$ SysAid $\cdot$ Alta}), senza che l'utente debba cambiare strumento.}
  \label{fig:orion_switch}
\end{figure}

\begin{figure}[htbp]
  \centering
  \includegraphics[width=0.95\textwidth,keepaspectratio]{approfondimenti-routing}
  \caption[Orion: meccanismo di clarification]{Il meccanismo di \emph{clarification}: quando la confidence del router è sotto la soglia, Orion non instrada ma pone all'utente una domanda di disambiguazione (``Su cosa vorresti fare l'analisi?'', ``A quale aspetto di Bologna ti riferisci?'') prima di procedere.}
  \label{fig:orion_clarification}
\end{figure}

\subsection{Renderer adattivo per risposte eterogenee}
\label{subsec:orion-impl-fe-shapes}

Il frontend implementa un renderer che si adatta 
alla \emph{shape} della risposta in base al campo \texttt{agent\_used}:

\begin{itemize}
  \item \textbf{SysAid / AllyCare:} la risposta può includere grafici (\texttt{charts}: 
  array di immagini base64), tabelle dati (\texttt{data}: JSON) e 
  query SQL (\texttt{queries}: array di stringhe). I grafici vengono renderizzati in coda alla risposta, 
  i dati in una tabella espandibile con download CSV, le query in un blocco codice collassabile.
  \item \textbf{CooPolicy:} la risposta include citazioni (\texttt{sources}: array con titolo, snippet, URI normalizzata). 
  Le citazioni sono rese come link cliccabili che aprono il PDF originale attraverso il proxy \texttt{/api/orion/docs/file}. 
  Le tabelle Markdown eventualmente presenti nella risposta sono renderizzate con \texttt{remark-gfm}.
  \item \textbf{Self (Orion):} la risposta segue il formato di Vera, con citazioni web inline e supporto a documenti caricati.
\end{itemize}

\begin{figure}[htbp]
  \centering
  \includegraphics[width=0.65\textwidth,height=0.9\textheight,keepaspectratio]{routing3}
  \caption[Orion: instradamento a CooPolicy con citazioni]{Una richiesta su una policy aziendale instradata a CooPolicy attraverso Orion: la risposta conserva le citazioni numerate \texttt{[N]} ai documenti sorgente, normalizzate dal proxy e rese cliccabili dal frontend.}
  \label{fig:orion_coopolicy}
\end{figure}

% =====================================================================
\section{Admin Panel di Orion}
\label{sec:orion-impl-admin}
% =====================================================================

Il pannello di amministrazione di Orion è composto da un frontend React
(\texttt{orion-admin-frontend}) e da un backend FastAPI dedicato
(\texttt{orion-admin-backend}), accessibili dal portale amministrativo.
In continuità architetturale con Vera, che dispone della coppia
\texttt{vera-admin-backend} / \texttt{vera-admin-frontend},
il backend admin di Orion è un servizio Cloud Run separato dal backend
principale, derivato dalla stessa base di \texttt{vera-admin-backend} e
operante sulle tabelle \texttt{orion\_*} del database condiviso.
I suoi endpoint sono serviti sotto il prefisso \texttt{/api/orion-admin/}.

Il pannello condivide con quello di Vera le pagine di gestione standard: budget token per utente, statistiche di consumo, log conversazioni, feedback strutturato, 
gruppi di esenzione PII e impostazioni Model Armor. Inoltre aggiunge le sezioni specifiche per l'orchestratore:

\begin{description}
  \item[Routing Logs.] Tabella paginata delle decisioni di routing, con timestamp, username, estratto del prompt, 
  agente selezionato, complessità, confidence (con codifica cromatica: verde $\geq 0{,}80$, giallo $\geq 0{,}60$, rosso $< 0{,}60$), flag di clarification. Il click su una riga espande il \emph{reasoning} completo del router, consentendo all'amministratore di diagnosticare le classificazioni errate.

  \item[Router Prompt.] Editor del \emph{system prompt} del router con textarea a tutta larghezza, 
  indicazione della data dell'ultimo aggiornamento e bottone ``Salva e Attiva''. 
  Il prompt è persistito nella tabella \texttt{app\_settings} del database, non nel codice: 
  questo consente di iterare il prompt senza un nuovo rilascio del backend.

  \item[Metriche di Routing.] Dashboard con quattro KPI (richieste totali, confidence media, clarification rate, numero di agenti attivi) 
  e un grafico a barre della distribuzione delle richieste per agente.

  \item[Eval Routing.] Pagina di valutazione automatica del routing, descritta in dettaglio nella sezione successiva.
\end{description}

\subsection{Pagina di valutazione integrata}
\label{subsec:orion-impl-eval}

La pagina \emph{Eval Routing} (\texttt{/orion-admin/eval}) è il componente che consente di eseguire la valutazione del routing direttamente dal browser, 
senza necessità di script CLI o accesso alla rete corporate. La pagina è organizzata in tre tab:

\paragraph{Tab ``Test Suite''.} Mostra la tabella dei test case attivi (ID, prompt, agente atteso, complessità attesa, categoria) 
con un conteggio in testa e un bottone ``Importa CSV'' che consente di caricare il file \texttt{test\_suite\_orion.csv} dal proprio computer. 
L'import sovrascrive i test case esistenti.

\paragraph{Tab ``Lancia Eval''.} Contiene un campo per la versione del prompt (es. \texttt{v1.0}, \texttt{v1.1}) 
e un bottone ``Avvia Valutazione''. Al click, il frontend invia una richiesta \texttt{POST /api/orion-admin/eval/run} al backend, 
che per ciascun test case chiama la funzione \texttt{classify()} \emph{direttamente} (invocazione Python, nessuna chiamata HTTP), 
confronta l'output con l'atteso, salva il risultato nel database e ritorna il summary completo. Al termine, la pagina visualizza:

\begin{itemize}
  \item Quattro KPI card con codifica cromatica verde/rosso rispetto ai target: \emph{Agent Accuracy}, \emph{Complexity Accuracy}, \emph{Mean Confidence (correct)}, \emph{Mean Confidence (wrong)}.
  \item Breakdown per agente con barre orizzontali proporzionali all'accuracy.
  \item Lista dei casi errati, espandibile per mostrare il \emph{reasoning} del router su ciascun errore.
  \item Tabella completa dei risultati, collassabile, con icone di correttezza per agente e complessità.
\end{itemize}

L'invocazione diretta della funzione \texttt{classify()} (anziché una chiamata HTTP
all'endpoint \texttt{/classify} del backend utente) è resa possibile dalla presenza di
una copia del modulo \texttt{router\_agent.py} all'interno di \texttt{orion-admin-backend}:
la funzione è importata come modulo Python, evitando il \emph{round trip} HTTP e la
necessità di autenticazione per ciascun test case.

\paragraph{Tab ``Storico Run''.} Mostra la lista delle run passate con timestamp, versione del prompt e metriche aggregate. 
Il click su una run carica il dettaglio completo, consentendo il confronto tra versioni successive del prompt.

% =====================================================================
\section{Infrastruttura Terraform}
\label{sec:orion-impl-iac}
% =====================================================================

L'infrastruttura di Orion aggiunge quattro servizi Cloud Run al deployment Terraform esistente: \texttt{orion-backend}, \texttt{orion-frontend}, \texttt{orion-admin-backend} e \texttt{orion-admin-frontend}.
Ciascun servizio ha il proprio trigger Cloud Build collegato al rispettivo repository Bitbucket.
Le variabili d'ambiente specifiche di Orion (URL dei sub-agenti, soglia di confidence, modello del router) 
sono iniettate da Terraform attraverso il file \texttt{tfvars} dell'ambiente corrispondente, seguendo lo stesso pattern degli altri servizi della suite.

% =====================================================================
\section{Database e schema dei log}
\label{sec:orion-impl-db}
% =====================================================================

Il database di Orion è ospitato sulla stessa istanza Cloud SQL della suite (database \texttt{ai\_agents\_db}), 
con tabelle distinte identificate dal prefisso \texttt{orion\_}. Le tabelle principali sono:

\begin{itemize}
  \item \texttt{orion\_users}, \texttt{orion\_conversations}, \texttt{orion\_messages}: ereditate dalla struttura di Vera, gestiscono utenti, conversazioni e messaggi.
  \item \texttt{orion\_routing\_logs}: registra ogni decisione del router con username, estratto del prompt, agente selezionato, complessità, confidence, flag di clarification e reasoning completo.
  \item \texttt{orion\_eval\_test\_cases}: contiene i test case della suite di valutazione (ID, prompt, agente atteso, complessità attesa, categoria, note).
  \item \texttt{orion\_eval\_results}: contiene i risultati di ogni run di valutazione (run ID, versione del prompt, agente atteso vs prodotto, complessità attesa vs prodotta, confidence, reasoning, flag di correttezza).
\end{itemize}

Tutte le tabelle sono create automaticamente al primo avvio di ciascun servizio
(\texttt{orion-backend} e \texttt{orion-admin-backend}) tramite
\texttt{Base.metadata.create\_all()}, in modo idempotente: il primo servizio che
parte crea lo schema, il secondo lo trova già presente.
Il pattern è lo stesso descritto nel Capitolo~\ref{chap:agenti-implementazione}:
non si utilizzano framework di migrazione, e le eventuali evoluzioni dello schema
sono gestite con istruzioni \texttt{ALTER TABLE ... ADD COLUMN IF NOT EXISTS}
nel codice di inizializzazione.

\begin{figure}[htbp]
    \centering
    \includegraphics[height=0.9\textheight,keepaspectratio]{orion-er}
    \caption[Schema del database Orion]{Schema relazionale del database di Orion: utenti, conversazioni, messaggi, log di routing e tabelle di valutazione (colonne tratte dai modelli SQLAlchemy del backend).}
    \label{fig:orion_er}
\end{figure}
